\section{Перевірка базових теоретичних відомостей та демонстраційних прикладів}

\subsection{Теоретичний аналіз особливостей типів даних у мові Java}

Згідно з програмою досліджень, у межах лабораторної роботи було проведено детальний аналіз 
взаємозв’язку між типами змінних, примітивними та об'єктними структурами даних у мові 
Java. Було досліджено такі фундаментальні аспекти:

\begin{itemize}
    \item \textbf{Особливості цілочисельного ділення типів \texttt{int}:} 
    При діленні двох об'єктів типу \texttt{int} результат завжди округляється до цілого 
    числа у бік нуля, а дробова частина повністю відкидається. Для отримання точного 
    результату з плаваючою комою хоча б один з операндів повинен явно мати тип \texttt{double} 
    або \texttt{float}.
    \item \textbf{Похибка представлення чисел з плаваючою комою (\texttt{double}):} 
    Через специфіку двійкового представлення дробових чисел у пам'яті комп'ютера за 
    стандартом IEEE 754, деякі арифметичні операції можуть виконуватися з незначною 
    похибкою (наприклад, операція $0.1 + 0.2$ повертає значення $0.30000000000000004$). 
    Цей фактор є критичним при проектуванні точних обчислювальних систем.
    \item \textbf{Ізоляція та порівняння об'єктного типу \texttt{String}:} 
    На відміну від примітивних типів даних, змінні класу \texttt{String} зберігають не 
    саме значення, а посилання на область пам'яті об'єкта у кучі (\textit{Heap}). 
    Через це застосування оператора \texttt{==} порівнює виключно адреси самих посилань, а не 
    їхній зміст. Для посимвольного порівняння вмісту рядків обов'язковим є використання 
    вбудованого методу \texttt{equals()}.
\end{itemize}

\subsection{Перевірка виконання демонстраційних прикладів у терміналі}

Для перевірки вищезазначених теоретичних тверджень та виконання обов'язкового пункту №2 
завдання було запущено службовий метод\\
 \texttt{TheoryExamples.runAllExamples()}. 
На рисунку~\ref{fig:terminal_theory} представлено скриншот відповідної ділянки консольного 
виведення у терміналі розробки.

\begin{figure}[H]
    \centering
    \includegraphics[width=0.95\textwidth]{screenshots/theory_terminal.png}
    \caption{Ділянка консолі з результатами перевірки базових прикладів із методички лаб3.docx}
    \label{fig:terminal_theory}
\end{figure}

Як видно з результатів тестування (рис.~\ref{fig:terminal_theory}), всі системні логи компілятора 
та віртуальної машини Java повністю підтвердили особливості обробки пам'яті та примітивів, 
описані в теоретичній частині роботи.

\subsection{Лістинг службового класу \texttt{TheoryExamples.java}}

Нижче наведено повний вихідний код створеного компонента, що інкапсулює у собі демонстраційні 
алгоритми перевірки мовних конструкцій.

%\subsection{Клас TheoryExamples.java}
\inputjava{../src/ad251/gorbunova/TheoryExamples.java}